% Comandos sugeridos para compilar:
% latexmk -pdf main.tex
% ou:
% pdflatex main.tex
% bibtex main
% pdflatex main.tex
% pdflatex main.tex

\documentclass[12pt,oneside,a4paper,brazil]{abntex2}

\usepackage[utf8]{inputenc}
\usepackage[T1]{fontenc}
\usepackage{lmodern}
\usepackage{babel}
\usepackage{microtype}
\usepackage{indentfirst}
\usepackage{booktabs}
\usepackage{longtable}
\usepackage{amsmath}
\usepackage{listings}
\usepackage{xcolor}
\usepackage{url}
\usepackage{hyperref}
\usepackage[alf]{abntex2cite}

\setlength{\parindent}{1.25cm}
\setlength{\parskip}{0.2cm}
\OnehalfSpacing
\justifying

\lstdefinestyle{codigoC}{
    language=C,
    basicstyle=\ttfamily\small,
    keywordstyle=\color{blue},
    commentstyle=\color{gray},
    stringstyle=\color{red!60!black},
    numbers=left,
    numberstyle=\tiny,
    frame=single,
    breaklines=true,
    tabsize=4,
    captionpos=b
}

\titulo{Benchmark em C: comparação entre variáveis \texttt{long int} e \texttt{register long int} em Windows x64 e Linux x64}
\autor{Chrysthyan Matheus}
\local{Manaus/AM}
\data{2026}
\orientador{Prof. Joethe Moraes de Carvalho}
\instituicao{%
Ministério da Educação\\
Instituto Federal de Educação, Ciência e Tecnologia do Amazonas\\
Campus Manaus Centro\\
Curso de Ciência da Computação\\
Disciplina: Sistemas Digitais}
\preambulo{Relatório acadêmico apresentado ao Instituto Federal de Educação, Ciência e Tecnologia do Amazonas -- Campus Manaus Centro, como parte das atividades da disciplina Sistemas Digitais.}

\begin{document}

\imprimircapa

\begin{folhaderosto}
\begin{center}
{\ABNTEXchapterfont\large Ministério da Educação}\\
{\ABNTEXchapterfont\large Instituto Federal de Educação, Ciência e Tecnologia do Amazonas}\\
{\ABNTEXchapterfont\large Campus Manaus Centro}\\[1.5cm]
{\ABNTEXchapterfont\large Curso de Ciência da Computação}\\
{\ABNTEXchapterfont\large Disciplina: Sistemas Digitais}\\
{\ABNTEXchapterfont\large Professor: Joethe Moraes de Carvalho}\\[2.5cm]
{\ABNTEXchapterfont\bfseries\Large Benchmark em C: comparação entre variáveis \texttt{long int} e \texttt{register long int} em Windows x64 e Linux x64}\\[2.5cm]
\end{center}
\begin{flushright}
Aluno: Chrysthyan Matheus\\
E-mail: 202500895@ifam.edu.br\\
Turma/período: 2º Período\\
Data de entrega: 16/06/2026
\end{flushright}
\vfill
\begin{center}
Manaus/AM\\
2026
\end{center}
\end{folhaderosto}

\textual

\chapter{Introdução}

Benchmark é um procedimento utilizado para medir e comparar o desempenho de sistemas computacionais, programas, componentes de hardware ou ambientes de execução sob uma metodologia definida. Em trabalhos de avaliação de desempenho, a medição deve ser feita de forma controlada, com repetição das execuções e registro dos valores obtidos, pois fatores como carga do sistema operacional, frequência do processador, compilador e processos em segundo plano podem alterar os tempos observados \cite{jain1991art,hennessy2019computer}.

A medição de desempenho computacional é importante porque permite observar, com base em dados experimentais, como diferentes ambientes ou escolhas de implementação influenciam o tempo de execução de um programa. Neste relatório, o desempenho é analisado pela latência de processamento, isto é, pelo tempo necessário para concluir a execução de dois laços aninhados implementados em linguagem C.

O objetivo deste trabalho é comparar o tempo de execução de dois programas em C. O primeiro utiliza variáveis normais do tipo \texttt{long int} para os contadores dos laços; o segundo utiliza a palavra-chave \texttt{register} com variáveis \texttt{register long int}. Ambos executam dois laços aninhados de 0 até 60000 e medem o tempo com a função \texttt{clock()}, da biblioteca \texttt{time.h} \cite{kernighan1988c,cppreference_clock}.

O enunciado original da atividade mencionava a comparação entre uma plataforma Intel de 64 bits no Windows e uma plataforma ARM no Linux. Entretanto, neste experimento, os testes foram adaptados e realizados na mesma máquina em dual boot, comparando Windows x64 e Linux x64 no mesmo hardware. Portanto, este relatório não afirma que houve execução em ARM; a análise apresentada refere-se especificamente à comparação entre Windows 11 Intel 64 bits e Linux Intel 64 bits.

\chapter{Desenvolvimento}

\section{Ambiente de testes}

Os testes foram realizados em uma mesma máquina física com dual boot. Essa escolha reduz diferenças de hardware entre os ambientes e concentra a comparação no sistema operacional, no ambiente de compilação e nas condições de execução. As características informadas para o ambiente foram:

\begin{itemize}
    \item Processador: Intel Core i5-12400F;
    \item Memória RAM: 32 GB DDR5;
    \item Sistema operacional Windows: Windows 11, Intel 64 bits;
    \item Sistema operacional Linux: Linux, Intel 64 bits;
    \item Compilador no Windows: GCC, versão mais recente disponível no ambiente;
    \item Compilador no Linux: GCC, versão mais recente disponível no ambiente.
\end{itemize}

Como a distribuição e a versão exata do Linux e as versões numéricas do GCC não foram registradas nos arquivos do projeto, esses dados devem ser confirmados caso seja necessário maior rigor de reprodutibilidade.

\section{Programas avaliados}

Foram avaliados dois programas em linguagem C. O programa com variáveis normais declara os contadores como \texttt{long int x, y}. O trecho principal é apresentado a seguir.

\begin{lstlisting}[style=codigoC,caption={Trecho principal com variáveis long int}]
long int x, y;
clock_t begin = clock();

for (x = 0; x < 60000; x++)
    for (y = 0; y < 60000; y++)
        ;

clock_t end = clock();
\end{lstlisting}

O segundo programa utiliza \texttt{register long int x, y}. Em C, a palavra-chave \texttt{register} indica ao compilador que uma variável será muito usada e que, se possível, poderia ser armazenada em registrador. Em compiladores modernos, essa palavra-chave funciona apenas como uma sugestão, e o compilador pode ignorá-la conforme suas próprias estratégias de alocação de registradores \cite{kernighan1988c,gcc2026}. O trecho principal é apresentado a seguir.

\begin{lstlisting}[style=codigoC,caption={Trecho principal com variáveis register long int}]
register long int x, y;
clock_t begin = clock();

for (x = 0; x < 60000; x++)
    for (y = 0; y < 60000; y++)
        ;

clock_t end = clock();
\end{lstlisting}

A medição foi feita pela diferença entre os valores retornados por \texttt{clock()} antes e depois dos laços. O resultado foi convertido para segundos pela divisão por \texttt{CLOCKS\_PER\_SEC}.

\section{Metodologia de execução}

Cada programa foi executado 10 vezes em cada ambiente, resultando em quatro conjuntos de medições: Windows com \texttt{long int}, Windows com \texttt{register long int}, Linux com \texttt{long int} e Linux com \texttt{register long int}. A compilação foi realizada com \texttt{gcc -O0}, com o objetivo de evitar otimizações automáticas que pudessem remover ou alterar o comportamento dos laços vazios.

A média aritmética foi calculada pela soma dos tempos dividida pelo número de execuções. O desvio padrão foi calculado para indicar a dispersão dos tempos em relação à média. O speedup foi calculado conforme a Equação~\ref{eq:speedup}.

\begin{equation}
\label{eq:speedup}
Speedup = \frac{Tempo\ médio\ do\ sistema\ de\ referência}{Tempo\ médio\ do\ sistema\ comparado}
\end{equation}

Nesta análise, o Windows foi usado como sistema de referência para calcular o speedup do Linux, pois os tempos médios no Linux foram menores nos dois programas.

\section{Resultados obtidos}

\begin{table}[htbp]
\centering
\caption{Tempos de execução dos benchmarks em segundos}
\label{tab:resultados}
\begin{tabular}{ccccc}
\toprule
Execução & Windows \texttt{long int} & Windows \texttt{register} & Linux \texttt{long int} & Linux \texttt{register} \\
\midrule
1 & 1,187000 & 0,983000 & 1,052731 & 0,452373 \\
2 & 1,177000 & 0,974000 & 1,052259 & 0,453411 \\
3 & 1,205000 & 0,932000 & 1,052573 & 0,455823 \\
4 & 1,147000 & 0,975000 & 1,043363 & 0,452255 \\
5 & 1,208000 & 0,957000 & 1,055390 & 0,452102 \\
6 & 1,188000 & 0,953000 & 1,051428 & 0,452353 \\
7 & 1,204000 & 0,945000 & 1,045174 & 0,454206 \\
8 & 1,223000 & 0,955000 & 1,055350 & 0,455323 \\
9 & 1,197000 & 0,945000 & 1,050876 & 0,455767 \\
10 & 1,188000 & 0,961000 & 1,050271 & 0,454044 \\
\midrule
Média & 1,192400 & 0,958000 & 1,050941 & 0,453766 \\
Desvio padrão & 0,020711 & 0,015734 & 0,003919 & 0,001492 \\
\bottomrule
\end{tabular}
\end{table}

Com base nas médias, o speedup sem \texttt{register}, comparando Windows \texttt{long int} com Linux \texttt{long int}, foi:

\begin{equation}
Speedup_{sem\ register} = \frac{1,192400}{1,050941} \approx 1,13
\end{equation}

O speedup com \texttt{register}, comparando Windows \texttt{register long int} com Linux \texttt{register long int}, foi:

\begin{equation}
Speedup_{com\ register} = \frac{0,958000}{0,453766} \approx 2,11
\end{equation}

Também é possível observar a diferença interna entre os programas em cada sistema. No Windows, a média passou de 1,192400 s para 0,958000 s, o que indica uma redução de tempo no programa com \texttt{register}. No Linux, a média passou de 1,050941 s para 0,453766 s, indicando redução ainda maior no programa com \texttt{register}. Apesar disso, como \texttt{register} é apenas uma sugestão ao compilador, essa diferença deve ser interpretada com cuidado e limitada ao ambiente testado.

\chapter{Conclusão}

Os resultados obtidos indicaram que, no cenário analisado, o Linux Intel 64 bits apresentou melhor desempenho que o Windows 11 Intel 64 bits nos dois programas avaliados. Para o programa com variáveis \texttt{long int}, o tempo médio foi de 1,192400 s no Windows e 1,050941 s no Linux, resultando em speedup aproximado de 1,13 para o Linux. Para o programa com \texttt{register long int}, o tempo médio foi de 0,958000 s no Windows e 0,453766 s no Linux, resultando em speedup aproximado de 2,11 para o Linux.

O uso de \texttt{register} apresentou diferença relevante nos resultados coletados, principalmente no Linux. Entretanto, não se deve concluir que a palavra-chave sempre produzirá o mesmo ganho em outros programas, compiladores ou arquiteturas. Em compiladores modernos, a alocação de registradores é controlada pelo próprio compilador, e a palavra-chave \texttt{register} pode ser ignorada ou ter efeito dependente do contexto.

As diferenças observadas podem ter sido influenciadas por fatores como sistema operacional, implementação do compilador, processos em segundo plano, gerenciamento de recursos, frequência dinâmica do processador, temperatura e resolução da medição de tempo. Além disso, o benchmark executa laços vazios, o que representa um caso específico e limitado. Assim, os resultados não devem ser generalizados para todo o desempenho dos sistemas operacionais, mas apenas para o experimento realizado com os programas, compilação e máquina descritos neste relatório.

\postextual
\bibliography{referencias}

\end{document}
